\section{Introduction}

\subsection{Background and Rationale}
Clinical trial matching is the process of aligning individual patients with trials for which they meet the precise inclusion and exclusion criteria. Historically, this matching has relied on manual review by clinical research coordinators, which is slow, labor-intensive, and prone to oversight \citep{smith2020manual}. Given the rapid expansion of trial protocols globally, manual matching cannot scale, resulting in up to 80\% of trials failing to meet their recruitment timelines \citep{jones2022recruitment}.

To address this, clinical investigators and data scientists have turned to Artificial Intelligence (AI). AI techniques, particularly Natural Language Processing (NLP), can extract key clinical concepts (such as diagnoses, lab values, and genomic markers) from unstructured clinical trial descriptions and patient electronic health records (EHRs) \citep{apex2024nlp}. While several proprietary systems (such as Nexa TrialMatch and Synapse Cohort) claim high accuracy, academic literature showcases a heterogeneous landscape of methodologies ranging from rule-based parsers to modern generative transformers and large language models (LLMs). 

A preliminary search of the Cochrane Database of Systematic Reviews, MEDLINE, and the JBI Evidence Synthesis database revealed no existing systematic reviews or scoping reviews on this specific topic. A scoping review is uniquely suited here as it will map the broad, interdisciplinary literature across computer science and clinical medicine, establishing a taxonomy of existing tools and identifying key research gaps in evaluation methods and ethnic or demographic bias.

\subsection{Objectives and Research Questions}
The objective of this scoping review is to systematically map and analyze the literature on AI applications in clinical trial matching. Specifically, we seek to answer the research questions outlined in Box~\ref{box:research_questions}.

\begin{tcolorbox}[protocolbox, title={Box 1: Primary Research Questions}, label=box:research_questions]
\begin{enumerate}[label=\textbf{RQ\arabic*}:, leftmargin=3.5em, itemsep=4pt, parsep=0pt, topsep=4pt, partopsep=0pt]
  \item What artificial intelligence methodologies (e.g., machine learning, deep learning, NLP, LLMs) are currently deployed for matching patients to clinical trials?
  \item What clinical domains (e.g., oncology, cardiology, rare diseases) are most frequently represented in AI-driven matching literature?
  \item How is patient data structured and integrated (e.g., EHR structured data, unstructured physician notes, genomic profiles, patient-reported outcomes)?
  \item What evaluation metrics (e.g., precision, recall, F1-score, time-to-match, screen-fail rates) are used to assess the performance and clinical utility of these systems?
  \item What challenges, limitations, and ethical considerations (e.g., algorithmic bias, data privacy, model explainability) are discussed in the literature?
\end{enumerate}
\end{tcolorbox}
